
\begin{abstract}
We construct a deterministic path-ensemble theory on a compact periodic
finite-horizon Sinai billiard with a genuinely nonconjugate moving scatterer.
The configuration is the triangular lattice Lorentz gas with one circular
scatterer per cell and radius in a compact subinterval of
$(\sqrt3/4,1/2)$.  The quotient table is compact, the horizon is uniformly
finite, and the collision map preserves the natural Liouville cosine measure.
In fixed angular boundary coordinates this measure is independent of the
radius.

The mechanical observable is the first lattice coordinate of the actual cell
displacement between successive collisions.  Its twisted Lorentz transfer
operator has a simple leading eigenvalue near the origin.  We derive a local
large-deviation principle, prove strict convexity from the nondegenerate Lorentz displacement
local-limit packet, and define the endogenous conjugate field by
$\theta_R(j)=(\Lambda_R')^{-1}(j)=I_R'(j)$.  A new two-sided central-window
Gibbs-conditioning theorem identifies current-conditioned Liouville paths with
the canonical driven equilibrium, with the endpoint eigenfunction bias
removed.  The driven kernel is only a conditional representation: every full
trajectory remains a deterministic orbit of the same specular billiard.

The radius family is not a similarity family.  Two normal period-two orbits
have incommensurate radius responses, and the mean free-flight time has an
explicit strictly negative Santal\'o derivative.  These results provide the
single physical platform imported by the remaining four papers.
\end{abstract}
\maketitle
\noindent\textbf{Platform identifier:} \texttt{TL1-LIO-v1}.

\tableofcontents

\section{The periodic triangular Lorentz family}\label{sec:geometry}
Let
\[
 e_1=(1,0),\qquad e_2=\left(\frac12,\frac{\sqrt3}{2}\right),
 \qquad \Lambda=\Z e_1\oplus\Z e_2,
\]
and let $\mathbb T_\Lambda^2=\R^2/\Lambda$.  Its area is
\[
 A_\Lambda=\abs{\det(e_1,e_2)}=\frac{\sqrt3}{2}.
\]
For
\[
 R\in\mathcal I:=\left[\frac9{20},\frac{47}{100}\right]
\]
define
\[
 Q_R=\mathbb T_\Lambda^2\setminus \overline B(0,R).
\]
The lifted table is the complement of the disks
$\overline B(\ell,R)$, $\ell\in\Lambda$.  A particle moves at unit speed and
reflects by the Euclidean equal-angle law.

\begin{theorem}[Uniform compact finite-horizon family]\label{thm:finite-horizon}
For every $R\in\mathcal I$, the disks are disjoint and the periodic Lorentz
gas has finite horizon.  The free-flight bound can be chosen uniformly in
$R$.
\end{theorem}
\begin{proof}
The shortest nonzero lattice vector has length one and
$2R\le 47/50<1$, so distinct disks are disjoint.  A collision-free infinite
line of rational lattice direction $v\in\Lambda$ would lie in a corridor
between adjacent rows of lattice points parallel to $v$.  Their separation is
$A_\Lambda/\abs v\le\sqrt3/2$.  Such a corridor requires
$2R<A_\Lambda/\abs v$, whereas
\[
 2R\ge\frac9{10}>\frac{\sqrt3}{2}.
\]
An irrational-direction line is dense on the torus and cannot avoid the open
disk.  Hence there is no infinite free line.  If free paths had unbounded
length at $R=9/20$, compactness of the space of oriented lines modulo the
lattice would yield an infinite free line in the limit.  Thus the smallest
radius has finite horizon; increasing the radius can only shorten flights.
\end{proof}

Let $M=\T\times(-\pi/2,\pi/2)$, where $\alpha\in\T=\R/(2\pi\Z)$ is angular
position on the circular boundary and $\varphi$ is the post-collision angle
from the inward normal.  We write $T_R:M\dashrightarrow M$ for the collision
map in these common coordinates and $\tau_R$ for the actual free-flight time.

\begin{proposition}[Common Liouville collision law]\label{prop:liouville}
For every $R\in\mathcal I$, $T_R$ preserves
\[
 d\mu(\alpha,\varphi)=\frac1{4\pi}\cos\varphi\,d\alpha\,d\varphi.
\]
Thus the normalized collision law is independent of $R$ in the common angular
coordinate.
\end{proposition}
\begin{proof}
The standard invariant collision element is
$\cos\varphi\,ds\,d\varphi$, with $ds=R\,d\alpha$.  Its total mass is
$2\abs{\partial Q_R}=4\pi R$.  Normalization cancels the factor $R$.
See \cite{ChernovMarkarian2006}.
\end{proof}

